\documentclass[11pt]{article}
\usepackage{manual_docs_style}
\externaldocument{build/run_graph_schema}[run_graph_schema.pdf]
\externaldocument{build/compute_metrics_stage}[compute_metrics_stage.pdf]

\begin{document}

\docTitle{Stage: Oracle Solution}
\phantomsection\label{docstart:oracle_solution}

The oracle stage searches for the best non-ground-truth explanation of an intervention trajectory.
It is used to test whether an accepted sample is distinguishable from alternative root-cause parameters.

The stage comprises:

\begin{itemize}
\item \code{workflows/metrics/parameter_optimisation_over_all.py}: orchestrator.
\item \code{workflows/metrics/parameter_optimisation.py}: per-run optimization.
\item \code{workflows/metrics/generate_trajectory_counterfactul_and_get_measurement.py}: counterfactual trajectory generation and measurement.
\end{itemize}

\section*{Orchestrator Invocation}

\begin{DocCode}
env/bin/python workflows/metrics/parameter_optimisation_over_all.py \
  --model MODEL [MODEL ...] \
  --policy POLICY_ID \
  [--overwrite] \
  [--matlab-workers N] \
  [--only-tsENV-questions]
\end{DocCode}

Optimization flags from \code{parameter_optimisation.py}, such as \code{--max-iter} and \code{--coarse-grid-points}, may be forwarded.

\section*{Inputs}

For each model, the oracle reads:

\begin{itemize}
\item \code{plans/<policy_id>/run_nodes.jsonl}: canonical recipes and run IDs.
\item \code{plans/<policy_id>/run_edges.jsonl}: baseline/intervention/time-zero relationships.
\item \code{metrics/<policy_id>/eligibility_metrics.jsonl}: eligible intervention candidates.
\item \code{experiment_config.json}: exposed parameters, allowed intervals, observable signals, and timing.
\item \code{noise_adder.py}: model-specific noise support.
\item \code{runs/<run_id>/simulink_model.mdl}: materialized Simulink model for the target run.
\item \code{runs/<run_id>/data.parquet}: target trajectory.
\end{itemize}

During migration, the stage may also read legacy \code{model_run_specs.json}; new code should use the resolved run graph.

\section*{Candidate Collection}

If \code{--only-tsENV-questions} is not passed, the orchestrator iterates over eligible intervention records from \code{eligibility_metrics.jsonl} whose family is eligible.

If \code{--only-tsENV-questions} is passed, the orchestrator iterates over the unique intervention runs selected in the exported question manifest.

For each selected intervention, the stage recovers the parent baseline recipe, changed parameter, intervention value, and intervention time from \code{run_nodes.jsonl} and \code{run_edges.jsonl}.

\section*{Per-Run Optimization}

A typical per-run invocation is:

\begin{DocCode}
env/bin/python workflows/metrics/parameter_optimisation.py \
  --model BallDrop \
  --policy benchmark_v1 \
  --run-id run_7fd9daa0055fb3f412ef5e0ab8d4c5ba \
  --max-iter 10 \
  --coarse-grid-points 9
\end{DocCode}

\code{RUN_ID} must identify an intervention run.
For each possible root-cause parameter, the script optimizes a post-intervention value that best fits the target trajectory.
Candidate simulations are diagnostic and are not registered as benchmark samples.

\section*{Objective}

The optimizer minimizes a normalized post-intervention trajectory loss.
Schematic objective:

\begin{DocCode}
raw_output = Simulator(candidate_parameter_value)
candidate_i = preprocessing(raw_output)
target_i = preprocessing(raw_target)
raw_loss_i = MSE(candidate_i, target_i)
baseline_loss_i = MSE(no_change_baseline_i, target_i)
normalized_loss_i = raw_loss_i / max(baseline_loss_i, 1e-6)
objective = mean(normalized_loss_i)
\end{DocCode}

The preprocessing is model-specific and signal-specific as defined in the paper appendix.

\section*{Output}

Per-run output is written under:

\begin{DocCode}
models/simulink/<MODEL>/runs/<run_id>/optimisation_results.json
\end{DocCode}

The orchestrator also writes:

\begin{DocCode}
models/simulink/<MODEL>/metrics/<policy_id>/optimisation_summary.json
\end{DocCode}

The summary records, for each intervention run, the ground-truth intervention, best alternative intervention, optimization losses, and whether the best alternative remains distinguishable from the target.

\section*{Counterfactual Generation}

\code{generate_trajectory_counterfactul_and_get_measurement.py} uses each \code{optimisation_results.json} file to simulate the best non-ground-truth counterfactual trajectory.
It writes a counterfactual parquet artifact and a detectability/distinguishability summary.

\section*{Validation Rules}

\begin{itemize}
\item The target \code{run_id} must be an intervention node in \code{run_nodes.jsonl}.
\item The target run must have a matched baseline edge.
\item Candidate parameters must come from \code{experiment_config.json::exposed_variables.parameters}.
\item Candidate values must remain inside the configured \code{allowed_intervals}.
\item The oracle must not modify \code{run_nodes.jsonl}, \code{run_edges.jsonl}, or \code{model_record.json}.
\end{itemize}

\end{document}
